% Generated draft: requires mathematical and source review.
\documentclass[11pt]{article}
\usepackage[margin=25mm]{geometry}
\usepackage{fontspec}
\setmainfont{Latin Modern Roman}
\usepackage{amsmath,amssymb,mathtools,unicode-math}
\setmathfont{Latin Modern Math}
\usepackage{enumitem,needspace,xurl,hyperref}
\setlength{\parindent}{0pt}
\setlength{\parskip}{5pt}
\setlength{\emergencystretch}{3em}
\newcommand{\N}{\mathbb{N}}
\newcommand{\Z}{\mathbb{Z}}
\newcommand{\Q}{\mathbb{Q}}
\newcommand{\R}{\mathbb{R}}
\newcommand{\C}{\mathbb{C}}
\newcommand{\F}{\mathbb{F}}
\newcommand{\A}{\mathbb{A}}
\newcommand{\PP}{\mathbb{P}}
\newcommand{\E}{\mathbb{E}}
\newcommand{\eps}{\varepsilon}
\newcommand{\dd}{\,\mathrm{d}}
\begin{document}
\section*{UnsolvedMath ID: 11000053}
\textbf{Source problem code:} catalogue entry
\subsection{Definitions and mathematical statement}
For every group $G$, $H_1(G;\Z)\cong G/[G,G]$ is its full abelianization, including torsion. Problem 5.6 asks to compute
\[
 H_1(\mathcal I_g(k);\Z)\qquad\text{for every integer }k\geq2.
\]
The notation $\mathcal I_g(k)$ is not defined in the supplied excerpt. In the commonly used Johnson-filtration convention for a closed surface, with $\pi=\pi_1(S_g)$ and $\gamma_1\pi=\pi$, $\gamma_{r+1}\pi=[\pi,\gamma_r\pi]$, one sets
\[
 \mathcal I_g(k)=\ker\bigl(\operatorname{Mod}_g\to
 \operatorname{Out}(\pi/\gamma_{k+1}\pi)\bigr).
\]
This is a conditional interpretation, not an identification established by the excerpt. The Johnson filtration, the lower central series of the Torelli group, and congruence subgroups are different constructions. The requested answer is integral abelianization, not just rational first homology.

The following extracted section on generating sets is context. The source reports a lower bound $(4g^3-g)/3$ on the number of generators for $\mathcal I_g$, giving $1330$ at $g=10$, attributed to abelian quotients and [Jo5]. Its use of ``rank'' needs verification: in the standard closed Torelli convention, rational abelianization and the number of generators contributed by torsion must be distinguished. The source contrasts Humphries's $2g+1$ Dehn twists and Wajnryb's two generators for the mapping class group, and Johnson's finite-generation proof [Jo2]. The extracted size of Johnson's set is printed $O(2g)$; an omitted exponent is possible and no silent correction is made. It asks how to name and manage these generators and announces a further fundamental problem whose statement is absent.
\subsection{Short English statement}
Compute the full integral abelianization of every subgroup $\mathcal I_g(k)$ for $k\geq2$, after recovering the precise filtration convention. The accompanying generator-count discussion has unresolved rank and extraction issues.
\subsection{Sources}
\begin{itemize}
\item Local registry: \url{data/unsolvedmath/problems.json}, permanent UnsolvedMath ID $11000053$, exact source code \texttt{catalogue entry}.
\item Farb, \emph{Problems on Mapping Class Groups and Related Topics}, Problem 5.6, PDF page 45: \url{https://www.math.uchicago.edu/~farb/papers/mcgbook.pdf}. Supplied access date: 2026-07-29; extraction: pdf-font-structure; page-disambiguated numbering among 329 headings.
\item Supplied references: Johnson [Jo5], corollary after Theorem 5; Johnson [Jo2]; Humphries; Wajnryb [Wa]. The bibliography entries and preceding filtration definition are not supplied.
\item The registry reports Johnson's computation of Torelli first homology and partial generalizations, describing these as level subgroups without a precise definition or additional references.
\end{itemize}
\end{document}
